\section{相关工作}

\subsection{水溶解度与多溶剂条件预测}

ESOL 以分子结构描述符估计水溶解度，适用于溶剂固定时的分子间比较\cite{delaney2004}。在多溶剂体系中，同一溶质的溶解度随溶剂和温度变化。Boobier 等将溶剂化能和相互作用等物化量用于回归，并比较随机森林、ExtraTrees 等学习器\cite{boobier2020}。Vermeire 等将热力学关系与溶剂化自由能和溶剂化焓等子模型结合，覆盖多种溶剂和温度\cite{vermeire2022}。本文的观测字段是两种分子结构、温度与实测溶解度，因此采用由 SMILES 计算的固定描述符和直接的条件回归。

\subsection{分子表示与小样本回归}

Ye 和 Ouyang 用分子指纹与温度建立多溶剂回归，并评价了未见溶质上的预测表现\cite{ye2021}。Attia 等比较了固定描述符模型和分子图消息传递模型，在留出溶质的设定中检验跨分子迁移\cite{fastsolv2025}。这些工作提供了输入表示和评价策略的比较依据。本文在基础描述符树上引入更完整的结构信息和同项描述符交互，并评价两类学习器的互补性。ExtraTrees 在当前训练数据上用随机化切分拟合非线性关系\cite{extratrees2006}，TabPFN 利用预训练表格模型对给定训练上下文生成回归预测\cite{tabpfn2025}。两者可在同一验证对象和目标尺度下进行比较。

\subsection{未见药物评价与模型选择}

同一对象的多条观测跨训练和验证集合会改变泛化评价的含义\cite{leakage2023}。水溶解度研究也指出，数据集重叠、实验来源和化学状态会影响评估\cite{solubility2024}。本文把不同条件下的同一药物母体归到同一折，并在误差汇总时给予各母体相同权重。嵌套交叉验证将配置选择与外层误差估计分开\cite{varma2006}；若再用外层结果筛选完整流程，评价还应说明这一层选择及其影响\cite{cawley2010}。

\subsection{预测解释与温度效应}

Gopichand 等用 SHAP 分析水溶解度和有机溶剂溶解度模型所使用的结构与条件信息\cite{gopichand2026}。树模型的局部归因可按描述符来源聚合\cite{treeshap2020}，而改变输入得到的响应曲线需要结合实际观测支持范围阅读\cite{apley2020}。本文通过分支输出与特征归因解释模型计算，再在实测数据中考察同一药物--溶剂组合内的温度关联。将关联解释为干预效应需要额外的实验与识别假设\cite{janzing2020,pearl2009}。
